% ============================================================================
%  10-do-hieu-suat — Đo hiệu suất của chính mình, nền sinh học, chương trình 12 tuần
%  Ngày tạo: 2026-09-09 | Sửa gần nhất: 2026-09-09 | Ôn gần nhất: chưa
%  Dependency: toàn bộ các chương trước. Mức độ: nên nắm.
% ============================================================================
\documentclass[../hieu-suat.tex]{subfiles}
\begin{document}
\chapter{Đo chính mình: chỉ số, nền sinh học, và một chương trình mười hai tuần}
\ptrtarget{10}\ptrtarget{10-do-hieu-suat}
\dependency{tất cả các chương trước; đặc biệt \ptr{02} (đếm vòng phản hồi), \ptr{05} (sổ dự đoán) và \ptr{08} (thí nghiệm trên chính mình phải chịu cùng chuẩn như thí nghiệm khác).}

\begin{tomtat}
Chương cuối ghép mọi thứ thành thứ chạy được. Nó bắt đầu bằng lý do phải đo thay vì cảm nhận --- có bằng chứng trực tiếp rằng cảm giác chủ quan về năng lực nhận thức tách rời khỏi hiệu suất thật khi thiếu ngủ kéo dài --- rồi đề xuất sáu chỉ số tự đo bám sát luận điểm gốc ``đếm vòng phản hồi, không đếm giờ'', và kết thúc bằng một chương trình mười hai tuần. Chương trình đó chưa có ai kiểm chứng; tôi trình bày nó đúng như một thí nghiệm cỡ mẫu một người, có dự đoán ghi trước, theo đúng chuẩn của chương \ptr{08}.
\end{tomtat}

\begin{khoigoi}
\h{Vì sao không thể dùng cảm giác ``hôm nay tôi làm việc hiệu quả'' làm thước đo?}
\h{Thiếu ngủ một chút mỗi đêm thì hại tới đâu, và tôi có tự nhận ra không?}
\h{Chỉ số nào cho biết tôi đang khá lên chứ không phải chỉ đang bận?}
\end{khoigoi}

\section{Vì sao phải đo: cảm giác không theo kịp hiệu suất}

Van Dongen và cộng sự \cite{vandongen2003} gán ngẫu nhiên người tham gia vào các mức ngủ 4, 6 hoặc 8 giờ mỗi đêm trong 14 đêm liên tiếp, trong điều kiện phòng thí nghiệm. Kết quả: hạn chế ngủ ở mức vừa phải cũng gây suy giảm chức năng thần kinh--hành vi đáng kể, và cái giá tích luỹ theo thời gian thay vì ổn định lại \nD{}. Điều đáng nhớ nhất với người làm nghiên cứu là phần \emph{tự nhận thức}: mức buồn ngủ tự đánh giá tăng lúc đầu rồi chững lại, trong khi hiệu suất khách quan tiếp tục xuống gần như tuyến tính --- nghĩa là người tham gia không nhận ra mức suy giảm của chính mình \nD{} (điểm về mức tự đánh giá này tôi đọc qua các nguồn thứ cấp mô tả bài gốc; chưa mở được toàn văn, nên phải kiểm lại trước khi trích vào một bản thảo).

Cùng một hình thái sai lệch đã xuất hiện hai lần trong bài học này: ảo giác năng lực khi đọc lại tài liệu \cite{karpicke2009} và niềm tin sai của kỳ thủ rằng họ đang tìm lời giải khác \cite{bilalic2008}. Ba kết quả độc lập, cùng một kết luận \nM{}: hệ thống tự quan sát của con người không đủ phân giải để dùng làm thiết bị đo, nên phải có sổ ghi bên ngoài.

Rasch và Born \cite{rasch2013} bổ sung mặt còn lại: giấc ngủ không chỉ là thời gian không làm việc mà là giai đoạn củng cố trí nhớ, trong đó các biểu diễn mới được kích hoạt lại và ổn định hoá, đặc biệt trong giấc ngủ sóng chậm \nT{}. Theo cách đó, cắt ngủ để học thêm là đánh đổi trực tiếp giữa số giờ mã hoá và chất lượng củng cố \nM{}.

\section{Chi phí chuyển việc}

Leroy \cite{leroy2009} đặt tên cho hiện tượng \emph{dư âm chú ý}: khi chuyển từ việc A sang việc B, hoạt động nhận thức về A vẫn tiếp tục và làm giảm hiệu suất ở B; hai thí nghiệm cho thấy người ta khó rời khỏi một việc \emph{chưa hoàn thành}, và --- điểm ngược trực giác --- áp lực thời gian buộc phải kết thúc việc A lại giúp chuyển sang việc B tốt hơn \nD{}.

Áp dụng \nM{}: cái đắt không phải số lần chuyển việc mà là số việc \emph{bỏ dở}. Kết thúc một đơn vị công việc ở trạng thái đóng --- ghi xong dòng nhật ký, viết xong câu tiếp theo phải làm gì --- rẻ hơn nhiều so với để nó lơ lửng.

\section{Sáu chỉ số tự đo}

Bám đúng luận điểm gốc của chương \ptr{02}: đếm vòng phản hồi, không đếm giờ. Đây là đề xuất của tôi \nM{}, không phải bộ chỉ số đã được kiểm chứng.

\begin{enumerate}
\item \textbf{Số vòng phản hồi đóng trong tuần} --- mỗi vòng là: có dự đoán viết trước, có phép đo, có kết luận ghi lại. Đây là chỉ số chính.
\item \textbf{Thời gian tới phản hồi đầu tiên} của mỗi ý tưởng mới: từ lúc nảy ra đến lúc có một phép đo hoặc một người đọc nói nó sai chỗ nào. Mục tiêu là kéo xuống hàng ngày, không phải hàng tháng.
\item \textbf{Tỉ lệ thí nghiệm có \code{hypothesis.md} viết trước} trên tổng số thí nghiệm đã chạy.
\item \textbf{Điểm hiệu chỉnh của sổ dự đoán}: trong các dự đoán ghi ở mức tin 70\%, bao nhiêu phần trăm thật sự đúng \cite{mellers2015}?
\item \textbf{Số lần bị người khác bác bỏ mỗi tháng} --- nếu bằng không thì không phải tôi đúng, mà là tôi chưa đưa ai xem đủ sớm \cite{mercier2011,dunbar1997} \nM{}.
\item \textbf{Số lượt viết lại ở tầng cấu trúc} (không tính sửa câu) trên mỗi bản thảo \cite{sommers1980}.
\end{enumerate}

Kèm hai chỉ số nền: giờ ngủ thật (ghi lại, không ước lượng) \cite{vandongen2003} và số việc bỏ dở cuối ngày \cite{leroy2009}.

\begin{cambay}
Chỉ số bị bỏ ra ngoài danh sách một cách có chủ ý: số giờ làm việc, số dòng mã, số paper đã đọc. Cả ba đều đo hoạt động chứ không đo thông tin sửa lỗi thu được, và cả ba đều tăng được mà không cần khá lên chút nào \nM{}.
\end{cambay}

\section{Chương trình mười hai tuần}

Mỗi giai đoạn thêm đúng một thứ, để nếu hiệu suất đổi thì còn quy trách nhiệm được --- chính là nguyên tắc ``đổi một biến'' của chương \ptr{08} áp cho bản thân \nM{}.

\begin{table}[H]\centering\footnotesize
\caption{Chương trình mười hai tuần. Cột cuối là chương chứa bằng chứng và chi tiết thủ tục.}
\begin{tabularx}{\linewidth}{c L{4.6cm} Y c}
\toprule
\textbf{Tuần} & \textbf{Thêm vào} & \textbf{Cách biết là đang làm đúng} & \textbf{Chương}\\
\midrule
1--2 & Chỉ đo, không đổi gì: sáu chỉ số + giờ ngủ & Có số nền để so; không được sửa cách làm việc trong hai tuần này & \ptr{10}\\
3--4 & Sổ dự đoán: mọi thí nghiệm có dự đoán kèm xác suất viết trước & Tỉ lệ thí nghiệm có dự đoán trước \(>80\%\) & \ptr{05}, \ptr{08}\\
5--6 & Nghi thức tự phản biện nửa giờ trước khi tin bất kỳ kết quả nào & Mỗi tuần bắt được ít nhất một giả thuyết cạnh tranh mà trước đó bỏ sót & \ptr{05}\\
7--8 & Quy trình viết lại sáu lượt, và đưa bản xấu cho người đọc sớm & Có ít nhất một lượt viết lại làm đổi cấu trúc, không chỉ đổi câu & \ptr{04}\\
9--10 & Đo nhiễu nền trước mọi tuyên bố; danh mục tám mục & Không còn tuyên bố nào dựa trên chênh lệch nhỏ hơn nhiễu nền & \ptr{08}\\
11--12 & Lịch ôn giãn dần cho cây tài liệu, ôn bằng truy hồi và trộn nhánh & Trả lời được bảng truy hồi của chương đã học tám tuần trước & \ptr{07}\\
\bottomrule
\end{tabularx}
\end{table}

\section{Chính chương trình này cũng là một giả thuyết chưa được kiểm}

Cần nói thẳng \nM{}: không có công trình nào trong khảo sát này đo hiệu quả của một chương trình luyện tập lên năng suất hay chất lượng nghiên cứu. Từng thành phần có bằng chứng riêng --- truy hồi, giãn cách, xen kẽ, xét mặt đối lập, hồi tưởng từ tương lai, danh mục tái lập --- nhưng \emph{tổ hợp} thì chưa. Vì vậy cách trung thực để chạy nó là coi nó như một thí nghiệm cỡ mẫu một người, và áp đúng chuẩn của chương \ptr{08}: viết trước dự đoán (``sau 12 tuần, số vòng phản hồi mỗi tuần tăng từ \(a\) lên \(b\)''), biết rằng cỡ mẫu một người không cho phép kết luận nhân quả, và ghi lại cả kết quả âm.

\begin{cannho}
Nếu chỉ giữ một thói quen từ toàn bộ bài học: \textbf{viết dự đoán kèm xác suất trước khi nhìn kết quả}. Nó là điều kiện cần cho mọi vòng phản hồi khác --- không có nó thì không có gì để sai, và không có gì để sai thì không có gì để học \nM{}.
\end{cannho}

\begin{traloi}
\tl{Vì cảm giác tách khỏi hiệu suất: khi bị hạn chế ngủ nhiều ngày, mức buồn ngủ tự đánh giá chững lại trong khi hiệu suất khách quan tiếp tục giảm \cite{vandongen2003}.}
\tl{Hại tích luỹ chứ không bão hoà, và phần lớn người ta không tự nhận ra \cite{vandongen2003}; ngoài ra ngủ còn là giai đoạn củng cố những gì vừa học \cite{rasch2013}.}
\tl{Số vòng phản hồi đóng mỗi tuần, thời gian tới phản hồi đầu tiên, và độ hiệu chỉnh của sổ dự đoán --- không phải số giờ hay số dòng mã \nM{}.}
\end{traloi}

\begin{tukiem}
\item Vì sao ``số giờ làm việc'' bị loại khỏi danh sách chỉ số, dù nó dễ đo nhất?
\item Kết quả nào cho thấy tự đánh giá về trạng thái nhận thức của mình là không đáng tin, và nó lặp lại ở mấy chỗ khác trong bài học này?
\item Dư âm chú ý: điều gì ngược trực giác trong kết quả của Leroy, và nó đổi cách bạn kết thúc một buổi làm việc thế nào?
\item Viết dự đoán cho chính chương trình 12 tuần của bạn, kèm xác suất, ngay bây giờ.
\end{tukiem}

\begin{onbai}
\ar[3]{Cảm giác \(\neq\) hiệu suất}{Buồn ngủ tự đánh giá chững lại trong khi hiệu suất tiếp tục giảm \cite{vandongen2003}.}
\ar[3]{Chỉ số chính}{Số vòng phản hồi đóng mỗi tuần (có dự đoán trước, có phép đo, có kết luận).}
\ar[2]{Thời gian tới phản hồi đầu tiên}{Từ ý tưởng tới phép đo hoặc người phản biện đầu tiên; kéo về hàng ngày.}
\ar[2]{Dư âm chú ý}{Việc \emph{bỏ dở} mới đắt, không phải việc chuyển; kết thúc ở trạng thái đóng \cite{leroy2009}.}
\ar[2]{Ngủ là một phần của học}{Củng cố trí nhớ diễn ra trong giấc ngủ, đặc biệt sóng chậm \cite{rasch2013}.}
\ar[3]{Thói quen quan trọng nhất}{Viết dự đoán kèm xác suất trước khi nhìn kết quả.}
\ar[1]{Giới hạn của chương trình}{Tổ hợp chưa ai kiểm; chạy như thí nghiệm n=1 có dự đoán ghi trước.}
\end{onbai}

\end{document}
